\documentclass{SOP5}\usepackage[]{graphicx}\usepackage[]{xcolor}
% maxwidth is the original width if it is less than linewidth
% otherwise use linewidth (to make sure the graphics do not exceed the margin)
\makeatletter
\def\maxwidth{ %
  \ifdim\Gin@nat@width>\linewidth
    \linewidth
  \else
    \Gin@nat@width
  \fi
}
\makeatother

\definecolor{fgcolor}{rgb}{0.345, 0.345, 0.345}
\newcommand{\hlnum}[1]{\textcolor[rgb]{0.686,0.059,0.569}{#1}}%
\newcommand{\hlsng}[1]{\textcolor[rgb]{0.192,0.494,0.8}{#1}}%
\newcommand{\hlcom}[1]{\textcolor[rgb]{0.678,0.584,0.686}{\textit{#1}}}%
\newcommand{\hlopt}[1]{\textcolor[rgb]{0,0,0}{#1}}%
\newcommand{\hldef}[1]{\textcolor[rgb]{0.345,0.345,0.345}{#1}}%
\newcommand{\hlkwa}[1]{\textcolor[rgb]{0.161,0.373,0.58}{\textbf{#1}}}%
\newcommand{\hlkwb}[1]{\textcolor[rgb]{0.69,0.353,0.396}{#1}}%
\newcommand{\hlkwc}[1]{\textcolor[rgb]{0.333,0.667,0.333}{#1}}%
\newcommand{\hlkwd}[1]{\textcolor[rgb]{0.737,0.353,0.396}{\textbf{#1}}}%
\let\hlipl\hlkwb

\usepackage{framed}
\makeatletter
\newenvironment{kframe}{%
 \def\at@end@of@kframe{}%
 \ifinner\ifhmode%
  \def\at@end@of@kframe{\end{minipage}}%
  \begin{minipage}{\columnwidth}%
 \fi\fi%
 \def\FrameCommand##1{\hskip\@totalleftmargin \hskip-\fboxsep
 \colorbox{shadecolor}{##1}\hskip-\fboxsep
     % There is no \\@totalrightmargin, so:
     \hskip-\linewidth \hskip-\@totalleftmargin \hskip\columnwidth}%
 \MakeFramed {\advance\hsize-\width
   \@totalleftmargin\z@ \linewidth\hsize
   \@setminipage}}%
 {\par\unskip\endMakeFramed%
 \at@end@of@kframe}
\makeatother

\definecolor{shadecolor}{rgb}{.97, .97, .97}
\definecolor{messagecolor}{rgb}{0, 0, 0}
\definecolor{warningcolor}{rgb}{1, 0, 1}
\definecolor{errorcolor}{rgb}{1, 0, 0}
\newenvironment{knitrout}{}{} % an empty environment to be redefined in TeX

\usepackage{alltt}

% ---------------------------------------------------------
% SOP Metadata
% ---------------------------------------------------------

\title{Calibration of pH Meter}
\author{Your Name}

\date{\today}

\SOPno{PH-001}

\approved{Dr. Jane Smith}

\ReviseDate{2026-05-10}

% Optional overrides
% \Department{Environmental Chemistry Laboratory}
% \SOPType{Laboratory SOP}
% \SOPLogo{../../assets/pomona.pdf}

% ---------------------------------------------------------
% Document
% ---------------------------------------------------------
\IfFileExists{upquote.sty}{\usepackage{upquote}}{}
\begin{document}

\maketitle

\tableofcontents

\newpage

% ---------------------------------------------------------
\section{Purpose}
% ---------------------------------------------------------

\NP
This SOP describes the procedure for calibrating laboratory pH meters prior to analytical measurements.

\NP
Calibration ensures measurement accuracy and consistency between laboratory personnel.

% ---------------------------------------------------------
\section{Scope}
% ---------------------------------------------------------

\NP
This procedure applies to all benchtop pH meters used within the laboratory.

% ---------------------------------------------------------
\section{Required Materials}
% ---------------------------------------------------------

\begin{itemize*}
    \item pH meter
    \item Calibration buffers (4, 7, and 10)
    \item Deionized water
    \item Kimwipes
\end{itemize*}

% ---------------------------------------------------------
\section{Procedure}
% ---------------------------------------------------------

\subsection{Initial Inspection}

\NP
Inspect the electrode for cracks, salt buildup, or air bubbles.

\NP
Verify the electrode storage solution is filled appropriately.

\subsection{Calibration}

\NP
Rinse the probe using deionized water.

\NP
Immerse the electrode into the pH 7 buffer solution.

\NP
Wait for the reading to stabilize before accepting calibration.

\NP
Repeat the process using pH 4 and pH 10 buffers.

% ---------------------------------------------------------
\section{Documentation}
% ---------------------------------------------------------

\NP
Record calibration date, operator initials, and slope values in the laboratory logbook.

% ---------------------------------------------------------
\section{References}
% ---------------------------------------------------------

\begin{enumerate*}
    \item Manufacturer operating manual
    \item EPA Method 150.1
\end{enumerate*}

\end{document}
